% cap01/1.4_objetivos.tex — Objetivos general y específicos.
\section{Objetivos}
\subsection{Objetivo General}
Diseñar, desarrollar y documentar con precisión académica una plataforma web SaaS \textit{multi-tenant} a escala comercial, denominada Democra. Su fin primordial es viabilizar la gestión integral y centralizada de los procesos operativos, la administración del capital humano, los recursos logísticos y la gobernanza estratégica de organizaciones de voluntariado peruanas. Esta construcción debe cimentarse en una arquitectura desacoplada moderna, provista de mecanismos transversales de seguridad perimetral, criptografía, auditoría profunda forense y metodologías de aseguramiento de la calidad; aplicando para todo el proceso la innovadora metodología de diseño \textit{Document-Driven SDLC} (DDS) orientada a inteligencias artificiales.

\subsection{Objetivos Específicos}
\begin{enumerate}[label=\alph*.]
    \item Implementar una arquitectura de software distribuida y orientada a servicios (por capas) que logre un desacoplamiento total: cliente interactivo tipo \textit{Single Page Application} basado en React y Fiber, una interfaz de aplicación programable (API REST) operada bajo el entorno no bloqueante Node.js, y un poderoso almacén transaccional sobre PostgreSQL.
    \item Definir y desplegar un modelo transaccional de base de datos multiesquema con un estricto aislamiento lógico por inquilino, orquestado dinámicamente por variables de sesión (\texttt{tenant\_id}) e impermeabilizado mediante restricciones estructurales de \textit{Row Level Security} (RLS).
    \item Programar y habilitar un modelo perimetral IAM (\textit{Identity and Access Management}) de control de acceso basado en roles (RBAC) con denegación explícita por defecto. Adicionalmente, inyectar un \textit{trigger} asíncrono de auditoría universal de base de datos y lógica de retención mediante borrado lógico blando (\textit{soft delete}).
    \item Establecer y validar empíricamente una robusta estrategia de aseguramiento de la calidad (QA). Esto abarca la implementación de tipado de datos estricto desde el inicio del proyecto, validación de \textit{inputs} en tiempo de ejecución y auditorías de seguridad continuas trazadas a lo largo de las fases de la metodología DDS.
    \item Validar con rigor científico y empírico el correcto funcionamiento de los dominios funcionales críticos de la plataforma y aplicar pruebas automatizadas de cobertura y pruebas unitarias para validar el correcto funcionamiento de los m\'odulos cr\'iticos.
\end{enumerate}
